\chapter{Conclusões e trabalhos futuros}
\label{ch:conclusoes}

\section{Conclusão principal}
\label{sec:conclusao-principal}

O \CabotageLens{} permite comparar uma viagem rodoviária direta com uma cadeia rodoviária--cabotagem--rodoviária. As duas alternativas mantêm a mesma origem, o mesmo destino e a mesma carga. O resultado mostra as pernas percorridas, as distâncias, os consumos, os custos modelados, as emissões operacionais e a fonte dos principais dados.

A principal contribuição é tornar o cálculo marítimo explicável. O sistema não supõe que o navio navegou somente entre os dois portos apresentados. Ele lê a sequência completa de escalas da ANTAQ, reconstrói a carga a bordo e calcula cada subtrecho entre a origem e o destino.

\section{Contribuição do modelo marítimo}
\label{sec:contribuicao-conclusao}

Uma viagem é elegível quando o porto de origem aparece antes do porto de destino. Isso inclui ligações diretas e viagens com qualquer número de escalas intermediárias. Para uma viagem Santos--Suape--Pecém--Manaus, por exemplo, o consumo observado da viagem é formado pelos três subtrechos: Santos--Suape, Suape--Pecém e Pecém--Manaus. Cada subtrecho usa sua distância e sua carga a bordo.

A intensidade é procurada primeiro pelo IMO no EU MRV. Quando não existe correspondência, o sistema usa uma estatística robusta da classe ou do tipo e registra esse fallback. O fallback resolve a ausência de um valor específico, mas continua sendo uma estimativa de grupo, e não uma medição do navio.

Todas as viagens elegíveis do mesmo OD participam da intensidade representativa. A regra usa a mediana ponderada pelo trabalho de transporte. O caminho do cenário é escolhido separadamente: ligação direta observada primeiro; na ausência dela, menor cadeia completa observada. Assim, diferentes corredores podem alimentar a intensidade sem serem unidos em uma rota artificial.

\section{Síntese das evidências}
\label{sec:alcance-evidencias-conclusao}

A matriz marítima contém 177 pares direcionais, 12.025 observações viagem--OD e 5.438 subtrechos. Entre os subtrechos, 3.290 usam intensidade por IMO e 2.148 usam fallback por tipo ou classe. Esses números informam simultaneamente a cobertura e a dependência de estimativas de grupo.

Santos--Manaus apresenta 89 observações em 22 corredores. Entre elas, 40 usam IMO e 49 usam fallback. Os fallbacks correspondem a 59,54\% do trabalho de transporte. A intensidade representativa é 9,322050~g/(t$\cdot$nm), enquanto a média aritmética ponderada diagnóstica é 25,243618~g/(t$\cdot$nm). A diferença mostra por que a mediana ponderada é menos sensível às intensidades altas.

O caminho selecionado é a ligação direta Santos--Manaus, com 6.112~km ou 3.300,216~nm. Os outros corredores, inclusive os que passam por Suape e Pecém ou por Rio de Janeiro e Salvador, permanecem na amostra de intensidade. Eles não alteram a geometria escolhida.

Para uma carga ilustrativa de 14~t, a navegação consome 430,71~kg de combustível segundo a intensidade e a distância selecionadas. Esse valor é apenas a parcela marítima do exemplo. Ele não representa a emissão porta a porta, o custo total ou uma cotação de frete.

O benchmark externo aponta a mesma direção modal em 21 pares positivos e suportados, mas com diferença relevante de magnitude. Portanto, ele oferece apoio direcional. Não oferece calibração quantitativa.

Em conjunto, os resultados mostram que o framework é útil para triagem acadêmica e comparação auditável. As emissões permanecem no escopo operacional \ttw{} de \cooe{}. Os custos permanecem valores modelados. Os resultados não demonstram superioridade universal da cabotagem, disponibilidade de serviço ou equivalência com inventários \wtw{}/LCA.

\section{Agenda prioritária de trabalhos futuros}
\label{sec:trabalhos-futuros}

\begin{table}[htbp]
\centering
\small
\caption{Agenda de extensão do \CabotageLens{}.}
\label{tab:agenda-futura}
\begin{tabularx}{\textwidth}{L{0.22\textwidth}Y Y}
\toprule
Eixo & Próximo passo & O que deve continuar registrado \\
\midrule
Dados marítimos & Ampliar a correspondência por IMO, revisar classes e incluir outros períodos ANTAQ/MRV. & Ano, fonte, cobertura, amostra, estatística e regra de outlier.\\
Incerteza & Calcular intervalos e testar a influência dos pesos e dos fallbacks. & Variabilidade entre viagens, incerteza da carga e incerteza do proxy MRV.\\
Rota e operação & Incluir serviço, frequência, janela, capacidade, terminal, tempo de trânsito, transbordo e confiabilidade. & Evidência específica de cada operador, porto e serviço.\\
Ambiental & Criar inventários separados para \wtw{}, LCA, combustíveis alternativos e emissões locais no porto. & Gases, energia, fatores, alocação e unidade funcional.\\
Econômico & Comparar o custo modelado com tarifas, margens, contratos, seguro e inventário. & Separação entre custo operacional, custo alocado e frete comercial.\\
Validação & Alinhar distâncias, carga, veículos, portos, alocação e fronteira com bases independentes. & Diferença entre direção, plausibilidade de magnitude e calibração.\\
\bottomrule
\end{tabularx}
\end{table}

Uma representação mais completa pode usar uma rede multimodal com linhas, operadores, frequências, capacidades e tempos de espera. Essa expansão depende de dados comerciais e operacionais que não podem ser inferidos apenas das viagens ANTAQ \cite{competitiveness2024,modalshiftreview2020}.

\section{Fechamento}
\label{sec:fechamento}

O \CabotageLens{} transforma a comparação modal em uma sequência verificável. O leitor consegue acompanhar a origem do dado, a ordem dos portos, a carga a bordo, a distância, a intensidade, o fallback e o critério estatístico.

Essa rastreabilidade é o que permite usar o resultado com responsabilidade. O sistema indica onde a cabotagem merece investigação e mostra quais premissas ainda precisam ser confirmadas antes de qualquer decisão operacional ou comercial.
